% !TEX program = xelatex
\documentclass[aspectratio=169,11pt]{beamer}
\usetheme{Madrid}
\usepackage{fontspec}
\setsansfont{DejaVu Sans}
\setmainfont{DejaVu Serif}
\usepackage{polyglossia}
\setdefaultlanguage{turkish}
\usepackage{amsmath,amssymb}
\usepackage{booktabs}
\usepackage{multicol}
\usepackage{xcolor}

% ---------- Renkler ----------
\definecolor{anamavi}{RGB}{0,45,98}
\definecolor{cozumyesil}{RGB}{0,110,50}
\definecolor{uyarikirmizi}{RGB}{180,20,20}
\setbeamercolor{structure}{fg=anamavi}
\setbeamercolor{frametitle}{bg=anamavi,fg=white}
\setbeamercolor{title}{bg=anamavi,fg=white}
\setbeamercolor{palette primary}{bg=anamavi,fg=white}
\setbeamercolor{palette secondary}{bg=anamavi!85!black,fg=white}
\setbeamercolor{palette tertiary}{bg=anamavi!70!black,fg=white}
\setbeamercolor{block title}{bg=anamavi,fg=white}
\setbeamercolor{block body}{bg=anamavi!7}
\setbeamercolor{block title example}{bg=cozumyesil,fg=white}
\setbeamercolor{block body example}{bg=cozumyesil!8}
\setbeamercolor{block title alerted}{bg=uyarikirmizi,fg=white}
\setbeamercolor{block body alerted}{bg=uyarikirmizi!7}
\setbeamertemplate{navigation symbols}{}
\setbeamertemplate{blocks}[rounded][shadow=false]

% ---------- Soru / Çözüm kalıbı ----------
% \soru{başlık}{soru metni}{çözüm metni}: önce soru, bir tıklamayla çözüm
\newcommand{\soru}[3]{%
\begin{frame}{#1}
  \begin{block}{Soru}\large #2\end{block}
  \pause
  \begin{exampleblock}{Çözüm}\large #3\end{exampleblock}
\end{frame}}
\newcommand{\ab}[1]{\overline{#1}}

\title[MMAT105 -- Hafta 3]{MMAT105 Genel Matematik}
\subtitle{Hafta 3: Bölünebilme Kuralları, EBOB ve EKOK}
\author{}
\date{2026--2027 Güz Dönemi}

\begin{document}

% =====================================================
\begin{frame}
\titlepage
\end{frame}

% (iletişim çerçevesi kaldırıldı)

\begin{frame}{Geçen Hafta Hatırlatma}
\begin{columns}[t]
\column{0.58\textwidth}
\begin{enumerate}
\item \(360\) sayısını asal çarpanlarına ayırınız.
\item \(360\) sayısının kaç tane pozitif tam böleni vardır?
\item Bir bölme işleminde bölen \(7\), bölüm \(12\), kalan \(5\) ise bölünen kaçtır?
\end{enumerate}
\column{0.38\textwidth}
\pause
\textbf{Cevaplar:}
\begin{enumerate}
\item \(360=2^3\cdot 3^2\cdot 5\)
\item \((3+1)(2+1)(1+1)=24\)
\item \(7\cdot 12+5=89\)
\end{enumerate}
\end{columns}
\vspace{6pt}
\begin{block}{Hatırlatma}
Bölünen \(=\) bölen \(\cdot\) bölüm \(+\) kalan, \qquad \(0\le\) kalan \(<\) bölen
\end{block}
\end{frame}

\begin{frame}{Bugün}
\begin{columns}[t]
\column{0.32\textwidth}
\begin{block}{Ders 1}
\begin{enumerate}
\item \(2,5,10\) kuralları
\item \(3\) ve \(9\) kuralı
\item \(4\) ve \(8\) kuralı
\item \(11\) kuralı
\item Kalan bulma, bilinmeyen rakam
\end{enumerate}
\end{block}
\column{0.32\textwidth}
\begin{block}{Ders 2}
\begin{enumerate}\setcounter{enumi}{5}
\item Bileşik bölünebilme: \(6,12,15,18,24,36,45,72\)
\item EBOB: tanım
\item Asal çarpan ve ortak bölme yöntemi
\end{enumerate}
\end{block}
\column{0.32\textwidth}
\begin{block}{Ders 3}
\begin{enumerate}\setcounter{enumi}{8}
\item EKOK
\item \(\text{EBOB}\cdot\text{EKOK}=a\cdot b\)
\item EBOB ve EKOK problemleri
\end{enumerate}
\end{block}
\end{columns}
\end{frame}

% =====================================================
\section{Ders 1: Bölünebilme Kuralları}
\begin{frame}
\centering\Large\color{anamavi}\textbf{Ders 1}\\[6pt]\normalsize Bölünebilme Kuralları ve Kalan Bulma
\end{frame}

\subsection{2, 5, 10}
\begin{frame}{\(2\), \(5\) ve \(10\) ile Bölünebilme}
\begin{block}{Kurallar (birler basamağına bakılır)}
\begin{itemize}
\item \(2\): Birler basamağı çift (\(0,2,4,6,8\)) ise bölünür. \ Kalan: birler basamağının \(2\)'ye bölümünden kalan.
\item \(5\): Birler basamağı \(0\) veya \(5\) ise bölünür. \ Kalan: birler basamağının \(5\)'e bölümünden kalan.
\item \(10\): Birler basamağı \(0\) ise bölünür. \ Kalan: birler basamağındaki rakam.
\end{itemize}
\end{block}
\begin{alertblock}{Sık yapılan hatalar}
\(10\) ile bölünen sayı hem \(2\) hem \(5\) ile bölünür; tersi de doğrudur (\(10=2\cdot 5\)).\\
\(5\) ile bölünüp \(2\) ile bölünmeyen sayının birler basamağı \(5\)'tir.
\end{alertblock}
\end{frame}

\soru{\(2\) ve \(5\) -- Örnek 1}
{\(\ab{47a}\) üç basamaklı sayısı \(2\) ile bölünebiliyor, ancak \(5\) ile bölünemiyor. \(a\) kaç farklı değer alabilir?}
{\(2\) ile bölünme: \(a\in\{0,2,4,6,8\}\).\\[3pt]
\(5\) ile bölünmeme: \(a\neq 0\) ve \(a\neq 5\).\\[3pt]
\(a\in\{2,4,6,8\}\) \(\Rightarrow\) \(4\) farklı değer.}

\subsection{3 ve 9}
\begin{frame}{\(3\) ve \(9\) ile Bölünebilme}
\begin{block}{Kural (rakamlar toplamına bakılır)}
\begin{itemize}
\item \(3\): Rakamları toplamı \(3\)'ün katı ise sayı \(3\) ile bölünür.
\item \(9\): Rakamları toplamı \(9\)'un katı ise sayı \(9\) ile bölünür.
\item \textbf{Kalan:} Sayının \(3\)'e (\(9\)'a) bölümünden kalan \(=\) rakamları toplamının \(3\)'e (\(9\)'a) bölümünden kalan.
\end{itemize}
\end{block}
\begin{alertblock}{Sık yapılan hatalar}
\(9\) ile bölünen her sayı \(3\) ile bölünür; tersi doğru değildir: \(12\) sayısı \(3\) ile bölünür, \(9\) ile bölünmez.\\
Bir rakam \(0\) da olabilir: \(a+9\) toplamı \(9\)'un katı ise \(a=0\) veya \(a=9\).
\end{alertblock}
\end{frame}

\soru{\(3\) ile Bölünebilme -- Örnek 2}
{\(\ab{5a72}\) dört basamaklı sayısı \(3\) ile bölünebiliyorsa \(a\)'nın alabileceği değerlerin toplamı kaçtır?}
{Rakamlar toplamı: \(5+a+7+2=14+a\).\\[3pt]
\(14+a\in\{15,18,21\}\) \(\Rightarrow\) \(a\in\{1,4,7\}\)\\[3pt]
Toplam: \(1+4+7=12\)}

\soru{\(9\) ile Bölünebilme -- Örnek 3}
{\(\ab{6a3}\) üç basamaklı sayısı \(9\) ile bölünebiliyorsa \(a\) hangi değerleri alabilir?}
{Rakamlar toplamı: \(6+a+3=9+a\).\\[3pt]
\(9+a\) sayısı \(9\)'un katı olmalı: \(9+a=9\) veya \(9+a=18\).\\[3pt]
\(a\in\{0,9\}\) \quad (Kontrol: \(603=9\cdot 67\), \ \(693=9\cdot 77\).)\\[3pt]
\textit{Dikkat:} \(a=0\) değerini unutmayın.}

\soru{Kalan Bulma -- Örnek 4}
{\(4\,738\,526\) sayısının \(9\) ile ve \(3\) ile bölümünden kalanlar kaçtır?}
{Rakamlar toplamı: \(4+7+3+8+5+2+6=35\)\\[3pt]
\(35=9\cdot 3+8\) \(\Rightarrow\) \(9\) ile bölümünden kalan \(8\).\\[3pt]
\(35=3\cdot 11+2\) \(\Rightarrow\) \(3\) ile bölümünden kalan \(2\).}

\soru{Kalan Bulma -- Örnek 5}
{\(\ab{2a7}\) üç basamaklı sayısının \(9\) ile bölümünden kalan \(4\) ise \(a\) kaçtır?}
{Rakamlar toplamı: \(2+a+7=9+a\).\\[3pt]
\(9+a\) sayısının \(9\) ile bölümünden kalan \(4\) olmalı: \(9+a=13\) \(\Rightarrow\) \(a=4\).\\[3pt]
(\(9+a=22\) için \(a=13\) rakam değildir.) \quad Kontrol: \(247=9\cdot 27+4\).}

\subsection{4 ve 8}
\begin{frame}{\(4\) ve \(8\) ile Bölünebilme}
\begin{block}{Kural (son basamaklara bakılır)}
\begin{itemize}
\item \(4\): Son iki basamağın oluşturduğu sayı \(4\)'ün katı ise (\(00\) dahil) sayı \(4\) ile bölünür.
\item \(8\): Son üç basamağın oluşturduğu sayı \(8\)'in katı ise (\(000\) dahil) sayı \(8\) ile bölünür.
\item \textbf{Kalan:} Sayının \(4\)'e (\(8\)'e) bölümünden kalan \(=\) son iki (üç) basamağın \(4\)'e (\(8\)'e) bölümünden kalan.
\end{itemize}
\end{block}
\begin{alertblock}{Sık yapılan hatalar}
Sadece birler basamağına bakmak yetmez: \(14\) çift ama \(4\)'e bölünmez; \(16\) ise bölünür.\\
\(8\) ile bölünen sayı \(4\) ile de bölünür; tersi her zaman doğru değildir.
\end{alertblock}
\end{frame}

\soru{\(4\) ile Bölünebilme -- Örnek 6}
{\(\ab{35a6}\) dört basamaklı sayısı \(4\) ile bölünebiliyorsa \(a\) kaç farklı değer alabilir?}
{Son iki basamak: \(\ab{a6}\). \ \(4\)'ün katı olan ve \(6\) ile biten iki basamaklı sayılar:
\[16,\ 36,\ 56,\ 76,\ 96\]
(\(06\) ise \(4\)'ün katı değildir.) \(a\in\{1,3,5,7,9\}\) \(\Rightarrow\) \(5\) farklı değer.}

\soru{\(4\) ve \(8\) ile Kalan -- Örnek 7}
{\(7\,253\,164\) sayısının \(4\) ile ve \(8\) ile bölümünden kalanlar kaçtır?}
{\(4\) için son iki basamak: \(64=4\cdot 16\) \(\Rightarrow\) kalan \(0\).\\[3pt]
\(8\) için son üç basamak: \(164=8\cdot 20+4\) \(\Rightarrow\) kalan \(4\).\\[3pt]
Sayı \(4\) ile bölünür, fakat \(8\) ile bölünmez.}

\subsection{11}
\begin{frame}{\(11\) ile Bölünebilme}
\begin{block}{Kural (sağdan sola \(+\,-\,+\,-\))}
Rakamlar, \textbf{birler basamağından başlayarak} sırasıyla \(+,-,+,-,\dots\) işaretleriyle toplanır.
\[\ab{abcde}\ \longrightarrow\ e-d+c-b+a\]
\begin{itemize}
\item Sonuç \(11\)'in katı ise (\(0\) ve negatif katlar dahil) sayı \(11\) ile bölünür.
\item \textbf{Kalan:} Sonucun \(11\)'e bölümünden kalandır; sonuç negatifse \(11\)'in katı eklenerek \(0\)--\(10\) arasına getirilir.
\end{itemize}
\end{block}
\begin{alertblock}{Sık yapılan hatalar}
İşaretler \textbf{sağdan} başlar. Soldan başlanırsa bölünebilme doğru çıkar, ama kalan yanlış çıkabilir.
\end{alertblock}
\end{frame}

\soru{\(11\) ile Bölünebilme -- Örnek 8}
{\(\ab{7a42}\) dört basamaklı sayısı \(11\) ile bölünebiliyorsa \(a\) kaçtır?}
{Sağdan: \(2-4+a-7=a-9\).\\[3pt]
\(a-9\) sayısı \(11\)'in katı olmalı; \(0\le a\le 9\) için tek olasılık \(a-9=0\).\\[3pt]
\(a=9\) \quad Kontrol: \(7942=11\cdot 722\).}

\soru{\(11\) ile Kalan -- Örnek 9}
{\(3052\) sayısının \(11\) ile bölümünden kalan kaçtır?}
{Sağdan: \(2-5+0-3=-6\).\\[3pt]
Sonuç negatif: \(-6+11=5\) \(\Rightarrow\) kalan \(5\).\\[3pt]
Kontrol: \(3052=11\cdot 277+5\).}

\subsection{Karma}
\soru{Bilinmeyen Rakam -- Örnek 10}
{\(\ab{3a5b}\) dört basamaklı sayısı \(5\) ile ve \(9\) ile bölünebiliyorsa \(a\)'nın alabileceği değerlerin toplamı kaçtır?}
{\(5\) ile bölünme: \(b\in\{0,5\}\). \ Rakamlar toplamı: \(8+a+b\).\\[3pt]
\(b=0\): \(8+a=9\) \(\Rightarrow\) \(a=1\) \quad (\(3150\))\\[3pt]
\(b=5\): \(13+a=18\) \(\Rightarrow\) \(a=5\) \quad (\(3555\))\\[3pt]
\(a\in\{1,5\}\) \(\Rightarrow\) toplam \(6\).}

\soru{Bilinmeyen Rakam -- Örnek 11}
{\(\ab{4ab}\) üç basamaklı sayısı \(10\) ile bölünebiliyor ve \(3\) ile bölümünden kalan \(2\)'dir. \(a\) kaç farklı değer alabilir?}
{\(10\) ile bölünme: \(b=0\). \ Rakamlar toplamı: \(4+a\).\\[3pt]
\(4+a\) sayısının \(3\) ile bölümünden kalan \(2\): \(4+a\in\{5,8,11\}\).\\[3pt]
\(a\in\{1,4,7\}\) \(\Rightarrow\) \(3\) farklı değer.}

\begin{frame}{Ders 1 -- Sınıfta Siz Çözün}
\begin{columns}[t]
\column{0.62\textwidth}
\begin{enumerate}
\item \(\ab{7a4}\) sayısı \(3\) ile bölünebiliyorsa \(a\) hangi değerleri alır? \(\star\)
\item \(58\,347\) sayısının \(9\) ile bölümünden kalan?
\item \(2\,358\,614\) sayısının \(8\) ile bölümünden kalan?
\item \(\ab{3a52}\) sayısı \(11\) ile bölünebiliyorsa \(a=?\) \(\star\)
\item \(4173\) sayısının \(11\) ile bölümünden kalan? \(\star\)
\end{enumerate}
\column{0.32\textwidth}
\pause
\textbf{Cevaplar:}
\begin{enumerate}
\item \(1,\ 4,\ 7\)
\item \(0\)
\item \(6\)
\item \(6\)
\item \(4\)
\end{enumerate}
\end{columns}
\vspace{4pt}{\footnotesize \(\star\) işaretli soruların adım adım çözümü sonraki slaytlardadır.}
\end{frame}

% ---- Ders 1 alıştırma çözümleri ----
\soru{Ders 1 -- Alıştırma Çözümü (Soru 1)}
{\(\ab{7a4}\) üç basamaklı sayısı \(3\) ile bölünebiliyorsa \(a\) hangi değerleri alır?}
{Rakamlar toplamı: \(7+a+4=11+a\); bu toplam \(3\)'ün katı olmalı.\\[3pt]
\(0\le a\le 9\) \(\Rightarrow\) \(11\le 11+a\le 20\) \(\Rightarrow\) \(11+a\in\{12,\ 15,\ 18\}\)\\[3pt]
\(a\in\{1,\ 4,\ 7\}\) \quad (\(714,\ 744,\ 774\))}

\soru{Ders 1 -- Alıştırma Çözümü (Soru 4)}
{\(\ab{3a52}\) dört basamaklı sayısı \(11\) ile bölünebiliyorsa \(a\) kaçtır?}
{Sağdan \(+,-,+,-\): \ \(2-5+a-3=a-6\).\\[3pt]
\(a-6\) sayısı \(11\)'in katı olmalı; \(0\le a\le 9\) için \(-6\le a-6\le 3\) \(\Rightarrow\) \(a-6=0\).\\[3pt]
\(a=6\) \quad Kontrol: \(3652=11\cdot 332\).}

\soru{Ders 1 -- Alıştırma Çözümü (Soru 5)}
{\(4173\) sayısının \(11\) ile bölümünden kalan kaçtır?}
{Sağdan \(+,-,+,-\): \ \(3-7+1-4=-7\).\\[3pt]
Sonuç negatif: \(11\) ekle \(\Rightarrow\) \(-7+11=4\).\\[3pt]
Kalan \(4\). \quad Kontrol: \(4173=11\cdot 379+4\).}

% =====================================================
\section{Ders 2: Bileşik Bölünebilme ve EBOB}
\begin{frame}
\centering\Large\color{anamavi}\textbf{Ders 2}\\[6pt]\normalsize Bileşik Bölünebilme ve EBOB
\end{frame}

\subsection{Bileşik bölünebilme}
\begin{frame}{Bileşik Sayılarla Bölünebilme}
\begin{block}{Kural}
Bölen, \textbf{aralarında asal} iki çarpana ayrılır; sayı her iki çarpanla da bölünmelidir.
\[6=2\cdot 3,\quad 12=3\cdot 4,\quad 15=3\cdot 5,\quad 18=2\cdot 9,\quad 24=3\cdot 8\]
\[36=4\cdot 9,\quad 45=5\cdot 9,\quad 72=8\cdot 9\]
(Aralarında asal: \(1\)'den başka ortak pozitif böleni olmayan sayılar.)
\end{block}
\begin{alertblock}{Sık yapılan hatalar}
\(12=2\cdot 6\) \textbf{kullanılamaz}: \(2\) ile \(6\) aralarında asal değildir.\\
Örnek: \(18\) sayısı \(2\) ile ve \(6\) ile bölünür, ama \(12\) ile bölünmez.
\end{alertblock}
\end{frame}

\soru{\(12\) Tuzağı -- Örnek 12}
{\(30\) sayısı \(2\) ile ve \(6\) ile bölünür. \(30\) sayısı \(12\) ile bölünür mü? Bir sayının \(12\) ile bölündüğünü göstermek için ne kontrol edilmelidir?}
{\(30=12\cdot 2+6\) \(\Rightarrow\) \(12\) ile \textbf{bölünmez}.\\[3pt]
\(2\) ve \(6\) aralarında asal değil (ortak böleni \(2\)); bu yüzden \(2\cdot 6\) ayrıştırması işe yaramaz.\\[3pt]
Doğru kontrol: \(12=3\cdot 4\) \(\Rightarrow\) hem \(3\) hem \(4\) ile bölünme.\\
\(30\) sayısı \(4\) ile bölünmediği için \(12\) ile de bölünmez.}

\soru{\(6\) ile Bölünebilme -- Örnek 13}
{\(\ab{5a4}\) üç basamaklı sayısı \(6\) ile bölünebiliyorsa \(a\)'nın alabileceği değerlerin toplamı kaçtır?}
{\(6=2\cdot 3\). Son rakam \(4\) çift \(\Rightarrow\) \(2\) ile her zaman bölünür.\\[3pt]
\(3\): \(5+a+4=9+a\) sayısı \(3\)'ün katı \(\Rightarrow\) \(a\in\{0,3,6,9\}\).\\[3pt]
Toplam: \(0+3+6+9=18\)}

\soru{\(6\) ile Kalan -- Örnek 14}
{\(\ab{2a5}\) üç basamaklı sayısının \(6\) ile bölümünden kalan \(1\) ise \(a\) kaç farklı değer alabilir?}
{Sayı \(6k+1\) biçimindedir: \(2\) ile bölümünden kalan \(1\) (sayı tek; son rakam \(5\), sağlanır) ve \(3\) ile bölümünden kalan \(1\).\\[3pt]
\(2+a+5=7+a\), \(3\) ile bölümünden kalan \(1\) \(\Rightarrow\) \(7+a\in\{7,10,13,16\}\)\\[3pt]
\(a\in\{0,3,6,9\}\) \(\Rightarrow\) \(4\) farklı değer. \quad (Kontrol: \(205=6\cdot 34+1\).)}

\soru{\(12\) ile Bölünebilme -- Örnek 15}
{\(\ab{41a8}\) dört basamaklı sayısı \(12\) ile bölünebiliyorsa \(a\) kaç farklı değer alabilir?}
{\(12=3\cdot 4\).\\[3pt]
\(4\): \(\ab{a8}\in\{08,28,48,68,88\}\) \(\Rightarrow\) \(a\in\{0,2,4,6,8\}\)\\[3pt]
\(3\): \(4+1+a+8=13+a\) \(\Rightarrow\) \(a\in\{2,5,8\}\)\\[3pt]
Ortak: \(a\in\{2,8\}\) \(\Rightarrow\) \(2\) farklı değer \quad (\(4128\), \ \(4188\)).}

\soru{\(15\) ile Bölünebilme -- Örnek 16}
{\(\ab{2a4b}\) dört basamaklı sayısı \(15\) ile bölünebiliyorsa \(a+b\) toplamı en çok kaçtır?}
{\(15=3\cdot 5\) \(\Rightarrow\) \(b\in\{0,5\}\). \ Rakamlar toplamı: \(6+a+b\).\\[3pt]
\(b=0\): \(6+a\) \(3\)'ün katı \(\Rightarrow\) \(a\in\{0,3,6,9\}\); \ en büyük \(a+b=9\)\\[3pt]
\(b=5\): \(11+a\) \(3\)'ün katı \(\Rightarrow\) \(a\in\{1,4,7\}\); \ en büyük \(a+b=12\)\\[3pt]
Cevap: \(12\) \quad (\(2745=15\cdot 183\)).}

\soru{\(45\) ile Bölünebilme -- Örnek 17}
{\(\ab{a37b}\) dört basamaklı sayısı \(45\) ile bölünebiliyorsa \(a\)'nın alabileceği değerler nelerdir?}
{\(45=5\cdot 9\) \(\Rightarrow\) \(b\in\{0,5\}\); \ \(a\neq 0\) (ilk rakam).\\[3pt]
\(b=0\): \(a+10\) \(9\)'un katı \(\Rightarrow\) \(a=8\) \quad (\(8370\))\\[3pt]
\(b=5\): \(a+15\) \(9\)'un katı \(\Rightarrow\) \(a=3\) \quad (\(3375\))\\[3pt]
\(a\in\{3,8\}\)}

\soru{\(36\) ile Bölünebilme -- Örnek 18}
{\(\ab{5a3b}\) dört basamaklı sayısı \(36\) ile bölünebiliyorsa bu koşulu sağlayan kaç farklı sayı vardır?}
{\(36=4\cdot 9\).\\[3pt]
\(4\): \(\ab{3b}\in\{32,36\}\) \(\Rightarrow\) \(b\in\{2,6\}\)\\[3pt]
\(9\): \(8+a+b\) \(9\)'un katı.\\
\(b=2\): \(a+10=18\Rightarrow a=8\); \quad \(b=6\): \(a+14=18\Rightarrow a=4\)\\[3pt]
Sayılar: \(5832\) ve \(5436\) \(\Rightarrow\) \(2\) farklı sayı.}

\soru{\(72\) ile Bölünebilme -- Örnek 19}
{\(\ab{a1b6}\) dört basamaklı sayısı \(72\) ile bölünebiliyorsa \(a+b\) kaçtır?}
{\(72=8\cdot 9\).\\[3pt]
\(8\): \(\ab{1b6}\) sayısı \(8\)'in katı \(\Rightarrow\) \(136\) veya \(176\) \(\Rightarrow\) \(b\in\{3,7\}\)\\[3pt]
\(9\): \(a+1+b+6=a+b+7\) \(9\)'un katı.\\
\(b=3\): \(a=8\); \quad \(b=7\): \(a=4\) \quad (\(8136\), \ \(4176\))\\[3pt]
Her iki durumda \(a+b=11\).}

\subsection{EBOB}
\begin{frame}{EBOB: En Büyük Ortak Bölen}
\begin{block}{Tanım ve yöntemler}
\(\text{EBOB}(a,b)\): \(a\) ve \(b\) pozitif tam sayılarının ortak bölenlerinin en büyüğü.
\begin{itemize}
\item \textbf{Asal çarpan yöntemi:} Sayılar asal çarpanlarına ayrılır; \textbf{ortak} asal çarpanlar \textbf{en küçük} üsleriyle çarpılır.
\item \textbf{Ortak bölme yöntemi:} Sayılar, hepsini birden bölen asal sayılarla bölünür; bölenlerin çarpımı EBOB'dur.
\end{itemize}
\end{block}
\begin{alertblock}{Dikkat}
EBOB, sayıların en küçüğünden büyük olamaz. \ Ortak asal çarpan yoksa EBOB \(=1\) (aralarında asal).
\end{alertblock}
\end{frame}

\soru{EBOB -- Örnek 20}
{\(\text{EBOB}(24,36)\) kaçtır?}
{\[24=2^3\cdot 3,\qquad 36=2^2\cdot 3^2\]
Ortak asal çarpanlar: \(2\) ve \(3\); en küçük üsler: \(2^2\) ve \(3^1\).
\[\text{EBOB}(24,36)=2^2\cdot 3=12\]
Kontrol: \(24=12\cdot 2\), \ \(36=12\cdot 3\) \ (\(2\) ve \(3\) aralarında asal).}

\soru{EBOB: Ortak Bölme -- Örnek 21}
{\(\text{EBOB}(84,126,210)\) kaçtır?}
{{\normalsize
\begin{tabular}{rrr|l}
\(84\) & \(126\) & \(210\) & \(2\)\\
\(42\) & \(63\) & \(105\) & \(3\)\\
\(14\) & \(21\) & \(35\) & \(7\)\\
\(2\) & \(3\) & \(5\) & \ (ortak bölen yok)
\end{tabular}}
\[\text{EBOB}=2\cdot 3\cdot 7=42\]}

\soru{EBOB: Üslü Biçim -- Örnek 22}
{\(A=2^3\cdot 3^2\cdot 5\) \ ve \ \(B=2^2\cdot 3^4\cdot 7\) \ olduğuna göre \(\text{EBOB}(A,B)\) kaçtır?}
{Ortak asal çarpanlar: \(2\) ve \(3\) \ (\(5\) ve \(7\) ortak değil).\\[3pt]
En küçük üsler: \(2^{\min(3,2)}=2^2\), \ \(3^{\min(2,4)}=3^2\).
\[\text{EBOB}(A,B)=2^2\cdot 3^2=36\]}

\begin{frame}{Ders 2 -- Sınıfta Siz Çözün}
\begin{columns}[t]
\column{0.62\textwidth}
\begin{enumerate}
\item \(\ab{8a6}\) sayısı \(6\) ile bölünebiliyorsa \(a\) kaç farklı değer alır?
\item \(\ab{a2b}\) üç basamaklı sayısı \(18\) ile bölünebiliyorsa \(a+b\) en çok kaçtır? \(\star\)
\item \(\text{EBOB}(60,90)=?\)
\item \(\text{EBOB}(48,72,120)=?\) \(\star\)
\item \(A=2^4\cdot 3\cdot 5^2\), \(B=2^2\cdot 3^3\cdot 7\) ise \(\text{EBOB}(A,B)=?\) \(\star\)
\end{enumerate}
\column{0.32\textwidth}
\pause
\textbf{Cevaplar:}
\begin{enumerate}
\item \(3\) \ (\(a\in\{1,4,7\}\))
\item \(16\) \ (\(828\))
\item \(30\)
\item \(24\)
\item \(2^2\cdot 3=12\)
\end{enumerate}
\end{columns}
\vspace{4pt}{\footnotesize \(\star\) işaretli soruların adım adım çözümü sonraki slaytlardadır.}
\end{frame}

% ---- Ders 2 alıştırma çözümleri ----
\soru{Ders 2 -- Alıştırma Çözümü (Soru 2)}
{\(\ab{a2b}\) üç basamaklı sayısı \(18\) ile bölünebiliyorsa \(a+b\) en çok kaçtır?}
{\(18=2\cdot 9\) (aralarında asal). \ \(2\): \(b\) çift. \ \(9\): \(a+2+b\) sayısı \(9\)'un katı.\\[3pt]
\(3\le a+b+2\le 20\) \(\Rightarrow\) \(a+b+2\in\{9,18\}\) \(\Rightarrow\) \(a+b\in\{7,16\}\).\\[3pt]
\(a+b=16\): \((a,b)=(8,8)\) uygun (\(b\) çift); \((7,9)\), \((9,7)\) uygun değil.\\[3pt]
En çok \(16\). \quad Kontrol: \(828=18\cdot 46\).}

\soru{Ders 2 -- Alıştırma Çözümü (Soru 4)}
{\(\text{EBOB}(48,72,120)\) kaçtır?}
{{\normalsize
\begin{tabular}{rrr|l}
\(48\) & \(72\) & \(120\) & \(2\)\\
\(24\) & \(36\) & \(60\) & \(2\)\\
\(12\) & \(18\) & \(30\) & \(2\)\\
\(6\) & \(9\) & \(15\) & \(3\)\\
\(2\) & \(3\) & \(5\) & \ (ortak bölen yok)
\end{tabular}}
\[\text{EBOB}=2\cdot 2\cdot 2\cdot 3=24\]}

\soru{Ders 2 -- Alıştırma Çözümü (Soru 5)}
{\(A=2^4\cdot 3\cdot 5^2\) \ ve \ \(B=2^2\cdot 3^3\cdot 7\) \ ise \(\text{EBOB}(A,B)\) kaçtır?}
{Ortak asal çarpanlar: \(2\) ve \(3\) \ (\(5\) yalnız \(A\)'da, \(7\) yalnız \(B\)'de).\\[3pt]
En küçük üsler: \(2^{\min(4,2)}=2^2\), \ \(3^{\min(1,3)}=3^1\).
\[\text{EBOB}(A,B)=2^2\cdot 3=12\]}

% =====================================================
\section{Ders 3: EKOK ve Problemler}
\begin{frame}
\centering\Large\color{anamavi}\textbf{Ders 3}\\[6pt]\normalsize EKOK, \(\text{EBOB}\cdot\text{EKOK}\) ve Problemler
\end{frame}

\subsection{EKOK}
\begin{frame}{EKOK: En Küçük Ortak Kat}
\begin{block}{Tanım ve yöntemler}
\(\text{EKOK}(a,b)\): \(a\) ve \(b\) pozitif tam sayılarının pozitif ortak katlarının en küçüğü.
\begin{itemize}
\item \textbf{Asal çarpan yöntemi:} \textbf{Bütün} asal çarpanlar \textbf{en büyük} üsleriyle çarpılır.
\item \textbf{Bölme tablosu:} Sayılar, en az birini bölen asal sayılarla, hepsi \(1\) olana kadar bölünür; bölenlerin çarpımı EKOK'tur.
\end{itemize}
\end{block}
\begin{alertblock}{Dikkat}
EKOK, sayıların en büyüğünden küçük olamaz. \ EBOB: ortak çarpanlar, küçük üs; \ EKOK: tüm çarpanlar, büyük üs.
\end{alertblock}
\end{frame}

\soru{EKOK -- Örnek 23}
{\(\text{EKOK}(12,18)\) kaçtır?}
{\[12=2^2\cdot 3,\qquad 18=2\cdot 3^2\]
Bütün asal çarpanlar en büyük üsleriyle: \(2^2\) ve \(3^2\).
\[\text{EKOK}(12,18)=2^2\cdot 3^2=36\]
Kontrol: \(36=12\cdot 3=18\cdot 2\).}

\soru{EKOK: Bölme Tablosu -- Örnek 24}
{\(\text{EKOK}(8,12,15)\) kaçtır?}
{\begin{columns}[c]
\column{0.45\textwidth}
{\normalsize
\begin{tabular}{rrr|l}
\(8\) & \(12\) & \(15\) & \(2\)\\
\(4\) & \(6\) & \(15\) & \(2\)\\
\(2\) & \(3\) & \(15\) & \(2\)\\
\(1\) & \(3\) & \(15\) & \(3\)\\
\(1\) & \(1\) & \(5\) & \(5\)\\
\(1\) & \(1\) & \(1\) &
\end{tabular}}
\column{0.5\textwidth}
\[\text{EKOK}=2^3\cdot 3\cdot 5=120\]
Bölünmeyen sayı aynen aşağı yazılır.
\end{columns}}

\subsection{EBOB·EKOK}
\begin{frame}{\(\text{EBOB}\cdot\text{EKOK}\) Bağıntısı}
\begin{block}{Kural (iki sayı için)}
\(a\) ve \(b\) pozitif tam sayılar ise
\[\text{EBOB}(a,b)\cdot\text{EKOK}(a,b)=a\cdot b\]
\(a\) ve \(b\) aralarında asal ise \(\text{EBOB}(a,b)=1\) ve \(\text{EKOK}(a,b)=a\cdot b\).
\end{block}
\begin{alertblock}{Sık yapılan hatalar}
Bu bağıntı \textbf{üç sayı için geçerli değildir}: \(2,4,8\) için \(\text{EBOB}\cdot\text{EKOK}=2\cdot 8=16\neq 64\).\\
Aralarında asal olmak için sayıların asal olması gerekmez: \(8\) ve \(15\) aralarında asaldır.
\end{alertblock}
\end{frame}

\soru{EBOB $\cdot$ EKOK -- Örnek 25}
{İki pozitif tam sayının çarpımı \(540\), EBOB'u \(6\) ise EKOK'u kaçtır?}
{\[6\cdot\text{EKOK}=540\ \Rightarrow\ \text{EKOK}=90\]
Örnek sayılar: \(6\) ve \(90\) \ (\(6\cdot 90=540\), \(\text{EBOB}=6\), \(\text{EKOK}=90\)).}

\soru{EBOB $\cdot$ EKOK -- Örnek 26}
{\(\text{EBOB}(a,b)=12\), \ \(\text{EKOK}(a,b)=180\) \ ve \(a=36\) ise \(b\) kaçtır?}
{\[36\cdot b=12\cdot 180=2160\ \Rightarrow\ b=60\]
Kontrol: \(36=2^2\cdot 3^2\), \(60=2^2\cdot 3\cdot 5\)\\[3pt]
\(\text{EBOB}=2^2\cdot 3=12\) \ ✓, \quad \(\text{EKOK}=2^2\cdot 3^2\cdot 5=180\) \ ✓}

\soru{Aralarında Asal -- Örnek 27}
{\(a\) ve \(b\) aralarında asal iki pozitif tam sayıdır ve \(\text{EKOK}(a,b)=72\)'dir. \(a+b\) toplamı en az kaçtır?}
{Aralarında asal \(\Rightarrow\) \(a\cdot b=\text{EKOK}=72=2^3\cdot 3^2\).\\[3pt]
\(2^3\) ve \(3^2\) bölünmeden ayrı sayılara gitmeli:
\[\{a,b\}=\{1,72\}\ \text{veya}\ \{8,9\}\]
En küçük toplam: \(8+9=17\)}

\subsection{EBOB problemleri}
\begin{frame}{EBOB mu, EKOK mu?}
\begin{columns}[t]
\column{0.48\textwidth}
\begin{block}{EBOB problemleri (parçalama)}
\begin{itemize}
\item Artmadan eşit parçalara kesme
\item Dikdörtgen alana en büyük kare fayans
\item Çevreye eşit aralıklı ağaç
\item ``En büyük'' ölçü, ``en az'' parça
\end{itemize}
\end{block}
\column{0.48\textwidth}
\begin{block}{EKOK problemleri (birleştirme)}
\begin{itemize}
\item Zillerin, otobüslerin yeniden birlikte çalışması
\item Hep aynı kalanı veren sayı
\item Küçük parçalardan en küçük kare
\end{itemize}
\end{block}
\end{columns}
\begin{alertblock}{Ağaç dikme}
\begin{itemize}
\item \textbf{Kapalı} çevre (bahçe, tarla): \ ağaç sayısı \(=\) çevre \(\div\) aralık
\item \textbf{Açık} yol (iki uca da ağaç): \ ağaç sayısı \(=\) uzunluk \(\div\) aralık \(+\,1\)
\end{itemize}
\end{alertblock}
\end{frame}

\soru{EBOB: Tel Kesme -- Örnek 28}
{Uzunlukları \(72\) cm, \(96\) cm ve \(120\) cm olan üç tel, hiç artmayacak biçimde eşit uzunlukta parçalara kesilecektir. En az kaç parça elde edilir?}
{En az parça için parça uzunluğu en büyük olmalı:
\[\text{EBOB}(72,96,120)=24\ \text{cm}\]
Parça sayısı: \(\dfrac{72}{24}+\dfrac{96}{24}+\dfrac{120}{24}=3+4+5=12\)}

\soru{EBOB: Kare Fayans -- Örnek 29}
{Boyutları \(180\) cm \(\times\) \(240\) cm olan dikdörtgen bir zemin, eş kare fayanslarla boşluksuz ve fayans kesilmeden döşenecektir. En az kaç fayans gerekir?}
{Fayans kenarı iki kenarı da tam bölmeli ve en büyük olmalı:
\[\text{EBOB}(180,240)=60\ \text{cm}\]
Fayans sayısı: \(\dfrac{180}{60}\cdot\dfrac{240}{60}=3\cdot 4=12\)}

\soru{EBOB: Ağaç Dikme -- Örnek 30}
{Kenarları \(90\) m ve \(150\) m olan dikdörtgen bir bahçenin çevresine, her köşede birer ağaç olacak ve ardışık ağaçlar arası uzaklıklar eşit olacak biçimde ağaç dikilecektir. En az kaç ağaç gerekir?}
{En az ağaç için aralık en büyük olmalı: \(\text{EBOB}(90,150)=30\) m.\\[3pt]
Çevre: \(2\cdot(90+150)=480\) m.\\[3pt]
Kapalı çevre: \(\dfrac{480}{30}=16\) ağaç \quad (\(+1\) \textbf{eklenmez}!)}

\subsection{EKOK problemleri}
\soru{EKOK: Ziller -- Örnek 31}
{Üç zil sırasıyla \(12\), \(18\) ve \(30\) dakikada bir çalmaktadır. Üçü birlikte saat 09{:}00'da çaldıktan sonra en erken saat kaçta yeniden birlikte çalarlar?}
{\[\text{EKOK}(12,18,30)=2^2\cdot 3^2\cdot 5=180\ \text{dakika}=3\ \text{saat}\]
\(09{:}00+3\ \text{saat}=12{:}00\)}

\soru{EKOK: Hep Aynı Kalan -- Örnek 32}
{Bir sınıftaki öğrenciler \(4\)'erli, \(5\)'erli ya da \(6\)'şarlı gruplara ayrıldığında her seferinde \(1\) öğrenci artmaktadır. Sınıf mevcudu \(1\)'den fazla ise en az kaçtır?}
{Mevcut \(x\) ise \(x-1\) sayısı \(4\), \(5\) ve \(6\)'nın ortak katıdır.\\[3pt]
\(\text{EKOK}(4,5,6)=60\) \(\Rightarrow\) \(x-1\in\{0,60,120,\dots\}\)\\[3pt]
\(x>1\) olduğundan en küçük \(x=60+1=61\).}

\soru{EKOK: Eksik Kalan -- Örnek 33}
{Bir pozitif tam sayı \(6\), \(8\) ve \(9\) ile bölündüğünde sırasıyla \(5\), \(7\) ve \(8\) kalanını veriyor. Bu koşulu sağlayan en küçük sayı kaçtır?}
{Her kalan, bölenden \(1\) eksik: sayıya \(1\) eklenirse üçüne de tam bölünür.\\[3pt]
\(x+1\) ortak kat \(\Rightarrow\) \(x+1=\text{EKOK}(6,8,9)=72\)
\[x=72-1=71\]
Kontrol: \(71=6\cdot 11+5=8\cdot 8+7=9\cdot 7+8\).}

\soru{EKOK: Üç Basamaklı -- Örnek 34}
{\(4\), \(6\) ve \(10\) ile bölündüğünde her seferinde \(3\) kalanını veren üç basamaklı en küçük doğal sayı kaçtır?}
{\(\text{EKOK}(4,6,10)=60\) \(\Rightarrow\) sayı \(60k+3\) biçiminde (\(k\in\mathbb{N}\)).\\[3pt]
\(k=1\): \(63\) (iki basamaklı) \quad \(k=2\): \(123\)\\[3pt]
Cevap: \(123\)}

\begin{frame}{Ders 3 -- Sınıfta Siz Çözün}
\begin{columns}[t]
\column{0.62\textwidth}
\begin{enumerate}
\item \(\text{EKOK}(18,24)=?\)
\item \(\text{EBOB}=4\), \(\text{EKOK}=120\), sayılardan biri \(24\) ise diğeri? \(\star\)
\item \(60\) cm ve \(84\) cm iki çubuk artmadan eşit parçalara kesilecek. En az kaç parça? \(\star\)
\item İki otobüs \(15\) ve \(20\) dakikada bir kalkıyor. 08{:}00'de birlikte kalktılarsa yeniden ne zaman?
\item \(3,4,5\) ile bölündüğünde hep \(2\) kalan üç basamaklı en küçük sayı? \(\star\)
\end{enumerate}
\column{0.32\textwidth}
\pause
\textbf{Cevaplar:}
\begin{enumerate}
\item \(72\)
\item \(20\)
\item \(12\) \ (\(\text{EBOB}=12\))
\item 09{:}00
\item \(122\)
\end{enumerate}
\end{columns}
\vspace{4pt}{\footnotesize \(\star\) işaretli soruların adım adım çözümü sonraki slaytlardadır.}
\end{frame}

% ---- Ders 3 alıştırma çözümleri ----
\soru{Ders 3 -- Alıştırma Çözümü (Soru 2)}
{İki pozitif tam sayının EBOB'u \(4\), EKOK'u \(120\)'dir. Sayılardan biri \(24\) ise diğeri kaçtır?}
{İki sayı için: \(\text{EBOB}\cdot\text{EKOK}=a\cdot b\).
\[24\cdot b=4\cdot 120=480\ \Rightarrow\ b=20\]
Kontrol: \(24=2^3\cdot 3\), \(20=2^2\cdot 5\) \(\Rightarrow\) \(\text{EBOB}=2^2=4\), \ \(\text{EKOK}=2^3\cdot 3\cdot 5=120\) \checkmark}

\soru{Ders 3 -- Alıştırma Çözümü (Soru 3)}
{\(60\) cm ve \(84\) cm uzunluğundaki iki çubuk, hiç artmayacak biçimde eşit uzunlukta parçalara kesilecektir. En az kaç parça elde edilir?}
{En az parça için parça boyu en büyük olmalı ve iki uzunluğu da tam bölmeli:\\[3pt]
\(60=2^2\cdot 3\cdot 5\), \ \(84=2^2\cdot 3\cdot 7\) \(\Rightarrow\) \(\text{EBOB}(60,84)=2^2\cdot 3=12\) cm\\[3pt]
Parça sayısı: \(\dfrac{60}{12}+\dfrac{84}{12}=5+7=12\)}

\soru{Ders 3 -- Alıştırma Çözümü (Soru 5)}
{\(3\), \(4\) ve \(5\) ile bölündüğünde her seferinde \(2\) kalanını veren üç basamaklı en küçük doğal sayı kaçtır?}
{Sayıdan \(2\) çıkarılırsa \(3\), \(4\) ve \(5\)'e tam bölünür.\\[3pt]
\(\text{EKOK}(3,4,5)=60\) \(\Rightarrow\) sayı \(60k+2\) biçiminde (\(k\in\mathbb{N}\)).\\[3pt]
\(k=1\): \(62\) (iki basamaklı) \quad \(k=2\): \(122\)\\[3pt]
Cevap: \(122\)}

\begin{frame}{Özet}
\begin{columns}[t]
\column{0.32\textwidth}
\begin{block}{Ders 1}
\small
\begin{itemize}
\item \(2,5,10\): birler basamağı
\item \(3,9\): rakamlar toplamı
\item \(4,8\): son iki / üç basamak
\item \(11\): sağdan \(+-+-\)
\item Kalan da aynı yoldan bulunur
\end{itemize}
\end{block}
\column{0.32\textwidth}
\begin{block}{Ders 2}
\small
\begin{itemize}
\item Bileşik bölen: aralarında asal çarpanlar (\(12=3\cdot 4\), \(2\cdot 6\) değil)
\item EBOB: ortak asal çarpanlar, en küçük üs
\end{itemize}
\end{block}
\column{0.32\textwidth}
\begin{block}{Ders 3}
\small
\begin{itemize}
\item EKOK: tüm asal çarpanlar, en büyük üs
\item \(\text{EBOB}\cdot\text{EKOK}=a\cdot b\) (iki sayı)
\item Kapalı çevre: çevre/aralık
\item Hep \(r\) kalan: \(\text{EKOK}\cdot k+r\)
\end{itemize}
\end{block}
\end{columns}
\vspace{6pt}
\begin{alertblock}{Gelecek hafta}
Rasyonel sayılar ve ondalık sayılar: dört işlem, sıralama
\end{alertblock}
\end{frame}

\begin{frame}
\centering\Large İyi çalışmalar
\end{frame}

\end{document}
